\Askhsh{23200}{4}
\begin{ylh}{1}
    \item 1.2 Συναρτήσεις
    \item 1.3 Μονότονες συναρτήσεις - Αντίστροφη συνάρτηση
    \item 2.6 Συνέπειες του Θεωρήματος Μέσης Τιμής.
\end{ylh}
Έστω $f:\mathbb{R}\to\mathbb{R}$ μια γνησίως μονότονη συνάρτηση της οποίας η γραφική παράσταση τέμνει τον άξονα $y'y$ στο σημείο με τεταγμένη 3 και διέρχεται από το σημείο $A(1, \ln2)$.
\begin{alist}
\item Να βρείτε τη μονοτονία της. \monades{5}
\item Να αποδείξετε ότι για οποιοδήποτε θετικό αριθμό $\alpha$ ισχύει $f(\alpha\ln\alpha) \le f(\ln\alpha)$. \monades{7}
\item Να λύσετε την εξίσωση $f(e^{x-1}+\ln x)=\ln2$. \monades{6}
\item Θεωρούμε τη συνάρτηση $g(x)=f(x)+(3-\ln2)x-3$, $x\in\mathbb{R}$. Να αιτιολογήσετε γιατί η συνάρτηση $g$ δεν αντιστρέφεται. \monades{7}
\end{alist}